% Se compila con XeLaTeX; hacen falta varias pasadas para que se resuelvan las referencias.
\documentclass[11pt,a4paper]{article}
\usepackage{fontspec}
\IfFontExistsTF{texgyrepagella-regular.otf}
 {\setmainfont{texgyrepagella}[Extension=.otf,UprightFont=*-regular,BoldFont=*-bold,
   ItalicFont=*-italic,BoldItalicFont=*-bolditalic]}
 {\IfFontExistsTF{Palatino Linotype}{\setmainfont{Palatino Linotype}}{\setmainfont{Latin Modern Roman}}}
\linespread{1.05}
% Con babel-spanish instalado (MiKTeX, TeX Live completo) se usa la opción clásica; si falta,
% se carga el locale español de babel sin patrones de guionado y los cortes de palabra se
% toman de guionado.tex.
\IfFileExists{spanish.ldf}
 {\usepackage[spanish,es-noshorthands,es-nodecimaldot,es-tabla]{babel}}
 {\usepackage{babel}\babelprovide[import,main,hyphenrules=nohyphenation]{spanish}% Puntos de corte de las palabras del artículo según la división silábica de la Ortografía de
% la lengua española (RAE y ASALE, 2010); no se separan vocales contiguas. Sólo se usa cuando
% la instalación de TeX no trae los patrones de guionado del español (véase preambulo.tex).
\babelhyphenation[spanish]{
ab-so-lu-ta ab-so-lu-ta-men-te ab-so-lu-tos ab-sor-be ab-sor-be-do-res ab-sor-be-dor
ab-sor-ber-lo aba-jo abier-to abier-tos ac-ti-va ac-ti-vas ac-ti-vi-dad ac-ti-vo ac-ti-vos
ac-túa ac-túan acer-ca aco-ta aco-ta-ción aco-ta-da aco-ta-das aco-ta-do aco-ta-dos
acu-mu-la-da acu-mu-la-do acu-mu-lan ad-mi-si-ble ad-mi-si-bles ad-mi-ta ad-mi-te ad-mi-ten
ad-mi-tir ad-ver-sa-ria-les ade-más afec-ta afir-ma afir-ma-cio-nes afir-ma-ción agre-ga-do
agru-pa-mos aho-ra ais-la-da ajus-tar al-can-ce al-can-za al-gún al-re-de-dor al-te-ra
al-ter-na-ti-va-men-te alea-to-ria alea-to-rio alea-to-rios ali-men-ta alu-de am-bas am-bos
an-chu-ra an-te-rio-res an-te-rior an-tes an-ti-dia-go-na-les an-ti-dia-go-nal ana-li-za
ana-lí-ti-ca ana-lí-ti-cas anu-la anu-lan anun-cia-da aná-li-sis aná-lo-gas apa-re-ce
apa-ri-ción apar-ta-dos apli-ca apli-ca-da apli-ca-do apli-ca-mos apli-car apli-car-la
apo-ya apor-tan apro-ve-chan apro-xi-ma apro-xi-ma-ción apén-di-ce ar-bi-tra-rias ar-chi-vo
ar-gu-men-to ar-ti-fi-cial ar-tí-cu-lo ar-xiv-per-fil aris-ta aris-tas ase-gu-ra asi-mis-mo
asig-na asig-na-do asig-na-mos asin-tó-ti-ca asin-tó-ti-ca-men-te asin-tó-ti-co aso-cia-dos
ata-ca ata-ca-do ata-can ata-que ata-ques au-men-ta au-men-ta-do au-xi-lia-res aun-que
avan-za axio-mas aña-de aña-den aña-dir ba-lan-ce bas-ta bas-ta-ría bi-la-te-ra-les
bi-na-ria bi-yec-ción blan-co blan-cos blo-que bo-rra bo-rra-do bo-rra-dos bo-rrar bor-de
bor-des bus-ca-da ca-bría ca-de-na ca-pas ca-rac-te-ri-za ca-re-ce ca-si-lla ca-si-llas
ca-sos cal-cu-la-do cam-bia cam-bian cam-biar cam-bio cam-bios can-ce-la can-ce-la-ción
can-ce-lan can-ce-lar-se can-di-da-ta can-di-da-tas car-te-sia-na car-te-sia-nas ce-ros
ce-rra-do cen-tros cer-ca cer-ca-no cer-ti-fi-ca-ción cer-ti-fi-ca-da cer-ti-fi-ca-do
cie-rre cin-co cir-cu-la-ri-dad cla-ve co-cien-te co-cien-tes co-de-gra-do co-lo-ca
co-lo-ca-cio-nes co-lo-ca-ción co-lo-ca-das co-lo-ca-do co-lo-ca-mos co-lo-can co-lo-car
co-lo-car-la co-lo-res co-lor co-lum-na co-lum-nas co-mu-nes co-mún co-nec-tan co-no-ce-mos
co-ro-la-rio co-rres-pon-de co-rres-pon-den co-rres-pon-dien-te co-tas co-te-jar
coe-fi-cien-te coe-fi-cien-tes coer-ci-vi-dad coin-ci-de com-bi-na-ción com-bi-na-to-ria
com-pa-ra com-pa-ra-ción com-pa-ran-do com-pa-rar com-pa-ti-ble com-pa-ti-bles com-pac-to
com-pac-tos com-par-ten com-pen-sa-ción com-pen-sa-dor com-pi-la-ción com-pi-la-dos
com-pi-lar-lo com-ple-men-ta-rio com-ple-ta com-ple-ta-bi-li-dad com-ple-ta-bles
com-ple-ta-cio-nes com-ple-ta-ción com-ple-ta-do com-ple-tan com-ple-tar com-ple-tar-se
com-ple-tas com-ple-ti-tud com-ple-to com-po-si-ción com-pro-ba-ción com-pro-bar com-prue-ba
con-ca-vi-dad con-cen-tra con-cen-tra-ción con-cen-tran con-cen-tre con-cier-ne
con-clu-sio-nes con-clu-sión con-di-cio-na con-di-cio-na-da con-di-cio-nal-men-te
con-di-cio-nar con-di-ción con-ge-la-mos con-ge-lan con-jun-ta con-jun-to con-mu-ta-cio-nes
con-se-cu-ti-vas con-se-cuen-cia con-se-cuen-cias con-ser-va con-ser-va-do-ra
con-ser-va-do-ras con-ser-van con-ser-ve con-si-de-rar con-si-de-re-mos con-si-guien-te
con-sis-te con-ta-do con-ta-do-res con-ta-dor con-ta-mos con-tar con-te-ni-da con-teo
con-teos con-ti-nua con-ti-nua-men-te con-ti-nuas con-ti-nui-dad con-ti-nuo con-tie-ne
con-tra-di-ce con-tri-bu-ción con-tro-la con-tro-la-da con-tro-lan con-tro-lar con-trol
con-ver-ge con-ver-gen con-ver-gen-cia con-vier-te cons-tan-te cons-tan-tes cons-tru-ye
cons-tru-yen cons-truc-ción cons-trui-da cons-trui-das cons-truir coor-de-na-da
coor-de-na-das cor-che-tes cor-ta cor-tan cor-te cre-ce cre-cer cre-cien-te cre-cien-tes
crea-do crez-ca cru-ce cru-ces cu-bier-to cu-bier-tos cu-bre cu-yas cu-yos cua-dra-do
cua-dra-tu-ra cua-dra-tu-ras cua-drá-ti-ca cua-drí-cu-la cua-ter-nas cua-tro cual-quie-ra
cual-quier cuan-do cuen-ta cuen-tas cues-ta cum-ple cum-plen có-di-go cón-ca-va cón-ca-vo
da-dos da-mos da-tos de-ci-ma-les de-ci-mal de-cir de-du-ce de-du-cir de-fi-ne de-fi-nen
de-fi-ni-cio-nes de-fi-ni-ción de-fi-ni-mos de-jan de-jar de-lan-te de-mos-tra-ción
de-mos-tra-da de-mos-tra-do de-mos-tra-mos de-mues-tra de-mues-tran de-más
de-no-mi-na-do-res de-no-mi-na-dor de-no-ta de-no-ta-mos de-no-tan de-no-te-mos de-pen-de
de-pen-den de-pen-den-cia de-pen-der de-re-cha de-ri-va de-ri-va-bi-li-dad de-ri-va-da
de-ri-va-das de-ri-van-do de-ri-vas de-sac-ti-va de-sac-ti-van de-si-gual-da-des
de-si-gual-dad de-sig-na de-te-ni-das de-te-ni-do de-ter-mi-na de-ter-mi-na-da
de-ter-mi-na-do de-ter-mi-nis-ta de-ter-mi-nis-tas de-ter-mi-nís-ti-ca-men-te de-tie-nen
den-si-da-des den-si-dad des-cri-be des-cri-ben des-cri-ta des-crip-cio-nes des-crip-ción
des-de des-ha-ce des-ha-cer des-pués des-ti-no des-ti-nos des-tru-ye des-truir
des-via-cio-nes di-fe-ren-cia di-fe-ren-cia-les di-fe-ren-cias di-fie-re di-fie-ren
di-men-sión di-ná-mi-ca di-re-mos di-rec-cio-nes di-rec-ción di-rec-ta di-rec-ta-men-te
di-rec-tas di-rec-to di-rec-tos di-ver-gen-cia di-vi-de di-vi-di-do di-vi-dien-do di-vi-dir
dia-go-na-les dia-go-nal dia-go-nal-men-te die-ci-séis dis-cre-ta dis-cre-tas
dis-cre-ti-za-ción dis-cre-to dis-cu-sión dis-po-ni-bi-li-dad dis-po-ni-ble dis-po-ni-bles
dis-tin-ción dis-tin-ta dis-tin-tas dis-tin-to dis-tin-tos dis-tri-bu-ción dis-tri-bu-ye
dis-tri-bui-da do-ble do-bles do-mi-nio do-mi-nios don-de du-ran-te dua-li-dad dé-bil
dé-bil-men-te dón-de ecua-cio-nes ecua-ción edi-to-rial efec-ti-vas efec-to ejem-plo
ejem-plos ela-bo-rar ele-gi-das ele-gi-do ele-gi-dos ele-gi-mos ele-gi-re-mos ele-gir
ele-men-tal ele-men-tos elec-cio-nes elec-ción eli-ge eli-gen eli-gien-do eli-mi-na
eli-mi-na-do eli-mi-na-dos eli-mi-nar em-pa-re-ja-mien-to em-pa-re-jan-do em-pal-me
em-pí-ri-ca en-cuen-tre en-la-ce en-te-ro en-te-ros en-ten-der en-ton-ces en-tre en-tro-pía
en-tró-pi-ca enu-me-ra-ción enun-cia enun-cia-do enun-cia-dos equi-va-le equi-va-len-cia
erro-res es-ca-la es-ca-lar es-ca-lo-na-da es-ca-lo-na-do es-ca-lo-na-dos es-co-ge es-cri-be
es-cri-ben es-cri-bi-mos es-cri-bien-do es-pa-cial es-pa-cio es-pe-ra-do es-pe-ran-za
es-pe-ran-zas es-pí-ri-tu es-ta-ble-ce es-ta-cio-na-rie-dad es-ta-do es-ta-rá es-ta-rán
es-tas es-ti-ma-cio-nes es-ti-ma-ción es-ti-mar es-tos es-truc-tu-ra es-tu-dia es-tu-dia-mos
es-tu-dian es-tu-diar es-tán es-tén eva-lua-da eva-lua-do eva-luar eva-lúa even-to even-tos
evi-ta evo-lu-ción ex-ce-den ex-cep-cio-na-les ex-cep-cio-nes ex-cep-ción
ex-clu-si-va-men-te ex-clu-sio-nes ex-clu-sión ex-clu-yen ex-clui-das ex-clui-dos
ex-haus-ti-va ex-pe-ri-men-to ex-pli-ca ex-plí-ci-ta ex-plí-ci-tas ex-plí-ci-to
ex-plí-ci-tos ex-po-nen-cial ex-pre-sa ex-pre-sar ex-pre-sión ex-ten-di-da ex-ten-sio-nes
ex-ten-sión ex-tien-de ex-tien-den ex-tra-po-la-ción ex-trac-tos ex-tre-mal ex-tre-mas
ex-tre-mo exac-ta exac-ta-men-te exac-tos exi-ge exi-gi-ría exi-to-sa exis-te exis-ten
exis-ten-cia exis-ten-cial fa-lla fa-llo fa-mi-lia fa-mi-lias fac-to-res fac-to-ri-za
fac-to-ri-za-ción fac-tor fal-ta fal-tan fi-bra fi-bras fi-gu-ra fi-gu-ras fi-ja-da
fi-ja-mos fi-jan fi-jan-do fi-jas fi-je-mos fi-jos fi-las fi-na-les fi-nal fi-nal-men-te
fi-ni-to fi-ni-tud fic-ti-cio fic-ti-cios fil-tro fle-chas for-ma for-ma-li-za-ción
for-ma-li-za-do for-ma-li-za-dos for-mal for-man frac-ción fran-ja fre-cuen-cia fren-te
fue-ra fuen-te fuen-tes fun-cio-na fun-cio-nal fun-cio-nes fun-ción fí-si-ca-men-te
fí-si-cas fí-si-co fí-si-cos fór-mu-la ge-ne-ra ge-ne-ra-do ge-ne-ra-dor ge-ne-ra-les
geo-me-tría geo-mé-tri-ca gra-dien-te gra-do gra-dos gran-de gran-des ha-bi-tual ha-ble-mos
ha-bía ha-ce-mos ha-cia ha-llar has-ta he-mos he-re-da he-rra-mien-ta he-rra-mien-tas
hes-sia-no heu-rís-ti-ca heu-rís-ti-ca-men-te hi-per-gra-fo hi-pó-te-sis ho-ri-zon-te
ho-ri-zon-tes hue-co hue-cos iden-ti-dad iden-ti-fi-ca iden-ti-fi-ca-ción iden-ti-fi-ca-mos
iden-ti-fi-can iden-ti-fi-car idén-ti-ca igua-les igual-dad ili-mi-ta-da ilus-tra-ción
im-pli-ca im-po-si-ble ima-gen in-ci-de in-clu-sión in-clu-so in-clu-ye in-clui-da
in-clui-do in-clui-dos in-clui-mos in-cre-men-tos in-de-pen-den-cia in-de-pen-dien-te
in-de-pen-dien-te-men-te in-de-pen-dien-tes in-de-xa-mos in-de-xan in-di-ca in-di-ca-ción
in-di-ca-do-res in-di-ca-dor in-di-can in-di-car in-di-vi-dua-les in-du-cen in-du-ci-das
in-fe-rior in-fe-rior-men-te in-fi-ni-ta in-for-ma-ción in-for-me in-me-dia-to
in-ser-cio-nes in-ser-ción in-ser-ta in-te-gra in-te-gra-bi-li-dad in-te-gra-do
in-te-gra-les in-te-gral in-te-gran in-te-gran-do in-te-grar in-te-li-gen-cia in-te-rior
in-ten-tar in-ter-pre-ta-ción in-ter-sec-ción in-ter-va-lo in-ter-va-los in-ter-vie-ne
in-tro-duc-ción in-ver-sa in-ver-sas in-ver-so in-ver-sos in-ver-tir in-ver-tir-se
inac-ti-vo ini-cial ini-cio ins-tan-te iso-me-tría ite-ra-do ja-co-bia-no je-rar-quía
jun-tas jun-to lar-go le-ga-les le-gal le-jos le-mas le-van-tar le-yes li-bre li-bres
li-neal lis-ta lla-ma-mos lla-ma-re-mos lle-gan lle-no lle-va lle-van lo-ga-rít-mi-co
lo-ga-rít-mi-cos lon-gi-tud lu-gar lue-go lí-mi-te lí-mi-tes lí-nea lí-neas ló-gi-ca ma-llas
ma-ne-ra ma-sas ma-xi-ma-li-dad ma-xi-mi-za ma-xi-mi-za-dor man-tie-ne mar-ca mar-cas
mar-gi-na-les mar-gi-nal mar-tin-ga-la mar-tin-ga-las me-di-ble me-di-da me-dia me-dian-te
me-dias me-dio me-jo-ra me-nor me-nos miem-bro miem-bros mien-tras mis-ma mis-mas mis-mo
mis-mos mix-tos mo-de-lo mo-de-ra-dos mo-men-to mo-ver mo-vi-mien-to mue-re muer-te mues-tra
mues-tran mues-trear mul-ti-pli-can-do mul-ti-pli-car má-xi-mo mé-to-do mó-du-lo mó-du-los
múl-ti-plo na-ci-mien-to na-ci-mien-tos ne-ce-si-ta ne-ce-si-ta-mos ne-ga-ti-va
ne-ga-ti-vi-dad neu-tro nin-gu-na nin-gún no-ta-ción no-ve-dad nor-ma nor-ma-li-za-ción
nor-ma-li-za-do nor-ma-li-za-do-res nor-ma-li-za-dor nor-ma-li-za-dos nor-ma-li-zan nu-los
nu-me-ra-mos nu-mé-ri-ca nu-mé-ri-co nu-mé-ri-cos nue-va nue-vas nue-vo nue-vos nues-tra
nú-cleo nú-me-ro ob-je-ti-vo ob-ser-va ob-ser-va-bles ob-ser-va-ción ob-te-ne-mos ob-te-ner
ob-te-ni-das ob-te-ni-do ob-te-ni-dos ob-tie-ne ob-tie-nen ob-tu-vo ocu-pa ocu-pa-da
ocu-pa-das ocu-pa-do ocu-pa-rían ocu-pan ocu-par ocu-rre ofre-ce omi-ti-mos op-ti-mi-za-ción
ope-ra-cio-nes ope-ra-ción ope-ra-do-res ope-ra-dor or-de-na-do or-de-na-dos or-den
or-di-na-ria or-di-na-rias or-di-na-rio or-ga-ni-zar or-to-go-na-li-dad ori-gi-na-les
ori-gi-nal orien-ta-cio-nes pa-la-bra pa-la-bras pa-pel pa-ra-da pa-ra-me-tri-za-da
pa-ra-mé-tri-ca pa-re-jas pa-res pa-rá-me-tro pa-rá-me-tros pa-sa-mos pa-san pa-sar pa-sos
pa-tro-nes pa-trón par-cial par-te par-ti-cu-lar par-ti-mos par-ti-to pe-di-do pe-que-ña
pe-que-ño pe-que-ños pe-sos pen-dien-te pen-dien-tes per-der per-fec-to per-fi-les per-fil
per-ma-ne-ce per-ma-ne-cen per-mi-te per-mi-ti-rá per-mi-tien-do per-mu-ta-ción
per-mu-to-nes per-mu-tón pie-za pie-zas po-de-mos po-drían po-ne-mos po-si-ble po-si-bles
po-si-cio-nes po-si-ción po-si-ti-va po-si-ti-vas po-si-ti-vi-dad po-si-ti-vo po-si-ti-vos
pon-de-ra-da pon-de-ra-das pon-de-ra-dos por-que pre-ce-de pre-ci-sa pre-gun-ta
pre-li-mi-na-res pre-pa-rar pre-ten-de pre-via-men-te pre-vio pre-vios prea-sig-na pri-me-ra
pri-me-ro pri-mer prin-ci-pa-les prin-ci-pal pro-ba-bi-li-da-des pro-ba-bi-li-dad pro-ba-mos
pro-bar pro-bar-la pro-ble-ma pro-ble-mas pro-ce-de pro-ce-den pro-ce-sa-mien-to pro-ce-so
pro-ce-sos pro-du-ce pro-du-cir-se pro-duc-to pro-duc-tos pro-gra-ma pro-lon-ga-do
pro-lon-gan pro-me-dia pro-me-diar pro-me-dio pro-me-dios pro-pia pro-pie-da-des
pro-po-si-cio-nes pro-po-si-ción pro-por-cio-na pro-por-cio-nal pro-por-ción pro-yec-cio-nes
pro-yec-ción pro-yec-ta pro-yec-tar prue-ba prue-ban prue-bas pró-xi-ma pu-ra-men-te pue-da
pue-dan pue-de pue-den pun-to pun-tos pun-tual pér-di-da que-da que-dan que-dar que-de
que-den ra-ros raí-ces re-ci-be re-ci-bir re-cla-ma re-co-no-ce re-co-rre re-cu-pe-rar
re-cí-pro-ca-men-te re-cí-pro-co re-cí-pro-cos re-don-dea-dos re-du-ce re-fe-ren-cia
re-fe-ri-re-mos re-fie-ren re-gión re-gla re-gu-la-ri-dad re-la-cio-na re-la-cio-nar-lo
re-la-cio-nes re-la-ción re-la-ti-va re-le-van-te re-lo-jes re-mon-ta re-pa-ra-cio-nes
re-pa-ra-ción re-pa-rar re-po-si-to-rio re-pre-sen-ta-ción re-pre-sen-ta-mos re-pro-du-ce
re-quie-re re-quie-ren re-si-dual re-si-duos re-so-lu-ción re-sol-ve-mos re-sol-vien-do
re-su-me re-suel-va re-suel-ve re-sul-ta re-sul-ta-do re-sul-ta-dos re-sul-tan-te re-ti-ra
re-ti-ra-das re-ti-ra-mos re-ti-rar re-ti-rar-se re-ti-re-mos re-ve-la re-ve-la-da
re-ve-la-do re-ve-lan re-ve-lar re-ve-lar-la re-vi-sio-nes re-vi-sión rea-les rea-li-zar
rea-pa-re-ce rees-ca-la rees-ca-la-da rees-ca-la-do rees-ca-la-dos rees-ca-la-mos
rees-ca-lar rees-truc-tu-ra-ción rei-na rei-nas res-pec-ti-va-men-te res-pec-to res-pon-de
res-to res-tos res-tric-cio-nes reu-ti-li-zó reú-ne sa-be-mos sa-ber sa-li-da sa-tis-fa-ce
sa-tis-fa-cen sal-tos sal-vo se-gui-mos se-gui-re-mos se-guir se-gun-da se-gun-do se-gun-dos
se-gún se-lec-cio-nar se-lec-ción se-pa-ra-do se-ría sec-cio-nes sec-ción seg-men-to
seg-men-tos sen-ti-do sen-ti-dos ses-go si-gue si-guen si-guien-te si-guien-tes
si-mul-tá-nea-men-te si-mé-tri-ca si-mé-tri-co si-tio siem-pre sien-do sig-no sim-ple sir-ve
sis-te-ma sis-te-mas so-bre so-bre-vi-ven so-lu-cio-nes so-lu-ción so-por-te
su-bu-ni-for-mes su-bín-di-ce su-bín-di-ces su-ce-si-va-men-te su-ce-si-vas su-ce-sión
su-fi-cien-te su-fi-cien-te-men-te su-gie-re su-man su-man-do su-man-dos su-mar su-mas
su-pe-rio-res su-pe-rior su-per-gra-dien-te su-per-vi-ven-cia su-po-ner su-pre-mo su-yos
sus-ti-tu-cio-nes sus-ti-tu-ción sus-ti-tu-ye sus-ti-tu-yen sus-ti-tui-mos sus-ti-tuir
ta-bla ta-ble-ro ta-ble-ros ta-les ta-ma-ño ta-ma-ños ta-sas tam-bién tam-po-co tan-to
tan-tos te-les-co-par te-ne-mos te-ner tem-po-ra-les tem-po-ral tem-pra-no ten-den-cias
ten-der ten-dien-te ten-ga ten-gan teo-re-ma teo-re-mas ter-mi-nal tex-to tex-tua-les ti-ras
tie-ne tie-nen tiem-po tiem-pos tien-de to-can to-da-vía to-das to-dos to-le-ran-cia
to-le-ran-cias to-ma-mos to-man to-mar to-roi-dal to-ta-les to-tal tra-ba-ja-mos tra-ba-jo
tra-ba-jos tra-mo tra-ta tra-tar tra-yec-to-ria tra-yec-to-rias tran-si-cio-nes tran-si-ción
trans-fe-ren-cia trans-fe-ren-cias trans-fe-rir trans-fie-re trans-ver-sal tras-la-cio-nes
tras-la-ción tras-la-da tras-la-da-da tras-la-dar tri-ples trian-gu-lar tér-mi-no tér-mi-nos
tí-pi-ca tí-pi-co tí-tu-lo um-bral uni-ci-dad uni-da-des uni-for-me uni-for-me-men-te
uni-for-mes uni-for-mi-dad usa-das usa-mos usa-re-mos usa-rá usan-do uti-li-dad
uti-li-za-ron va-can-cia va-can-cias va-cía va-cías va-cío va-cíos va-len va-lo-res va-lor
va-ria-ble va-ria-cio-nal va-ria-cio-nes va-ria-ción va-rias va-ría ve-ces ve-lo-ci-dad
ve-ri-fi-ca ve-ri-fi-ca-ción ver-da-de-ra ver-da-de-ras ver-da-de-ro ver-sio-nes ver-sión
vi-ce-ver-sa vir-tua-les vir-tual vis-ta vo-raz vol-va-mos vol-ver vol-vien-do vuel-ven
véa-se vér-ti-ce vér-ti-ces éxi-to ín-di-ce óp-ti-mo óp-ti-mos ór-de-nes úl-ti-ma úl-ti-mas
úl-ti-mo úl-ti-mos úni-ca úni-ca-men-te úni-co úti-les
}
}
\usepackage{amsmath,amssymb,amsthm,mathtools}
\usepackage{unicode-math}
\IfFontExistsTF{texgyrepagella-math.otf}{\setmathfont{texgyrepagella-math.otf}}{\setmathfont{Latin Modern Math}}
\usepackage[margin=26mm]{geometry}
\IfFileExists{microtype.sty}{\usepackage{microtype}}{}
\IfFileExists{titlesec.sty}{\usepackage[nobottomtitles*]{titlesec}
  \renewcommand{\bottomtitlespace}{.12\textheight}}{}
\usepackage{booktabs,array}
\usepackage[shortlabels]{enumitem}
\usepackage{tikz}
\usetikzlibrary{arrows.meta,positioning}
\usepackage{xcolor}
\usepackage[colorlinks=true,linkcolor={blue!45!black},citecolor={green!35!black},urlcolor={blue!55!black}]{hyperref}

\AtBeginDocument{\renewcommand{\tablename}{Tabla}}
\setlength{\parskip}{3pt}
\setlength{\emergencystretch}{2em}
\allowdisplaybreaks[2]
\numberwithin{equation}{section}

\theoremstyle{plain}
\newtheorem{teorema}{Teorema}[section]
\newtheorem{proposicion}[teorema]{Proposición}
\newtheorem{lema}[teorema]{Lema}
\newtheorem{corolario}[teorema]{Corolario}
\theoremstyle{definition}
\newtheorem{definicion}[teorema]{Definición}
\theoremstyle{remark}
\newtheorem{observacion}[teorema]{Observación}

% Notación
\newcommand{\E}{\mathbb{E}}
\newcommand{\PP}{\mathsf{P}}
\newcommand{\QQ}{\mathsf{Q}}
\newcommand{\KL}[2]{D(#1\,\Vert\,#2)}
\newcommand{\eps}{\varepsilon}
\newcommand{\ind}{\mathbf{1}}
\newcommand{\Om}{\Omega}
\newcommand{\cL}{\mathcal{L}}
\newcommand{\oc}{\mathrm{oc}}
\newcommand{\di}{\mathrm{di}}
\newcommand{\inv}{\mathrm{inv}}
\newcommand{\vac}{\mathrm{vac}}
\newcommand{\col}{\mathrm{col}}
\DeclareMathOperator{\Lip}{Lip}
